\subsection{Repairing the ranking}\label{sec:repair-res}
Table~\ref{tab:repair} reports the leave-one-corpus-out repairs (per-target values in Table~\ref{tab:repair-targets}). Training on the six other corpora at once is the one change that helps ranking: held-out AUC rises to \LocoPooled{}, against \LocoMeanSingle{} for the average single source. This is close to the best single source chosen with hindsight (\LocoBestSingle), which cannot be identified without target labels, and far below the targets' own in-domain AUC of \LocoCeiling{}. Nothing we add on top of pooling helps on average. Balancing every (corpus, class) cell (\LocoBal, $\Delta$\LocoBalDelta, $p=\LocoBalP$) and masking collection signatures (\LocoMask, $\Delta$\LocoMaskDelta, $p=\LocoMaskP$) leave AUC unchanged on average; masking helps clearly only on ISOT (Table~\ref{tab:repair-targets}), the corpus whose signatures it was written for.

Methods that make the representation \emph{invariant to corpus identity} make ranking worse: deleting the 2{,}000 most corpus-predictive features gives \LocoDrop{} ($\Delta$\LocoDropDelta, $p=\LocoDropP$), CORAL \LocoCORAL{} ($\Delta$\LocoCORALDelta, $p=\LocoCORALP$) and INLP \LocoINLP{} ($\Delta$\LocoINLPDelta, $p=\LocoINLPP$). INLP does what it is designed to do, reducing linear corpus-identity accuracy from \INLPaccBefore{} to \INLPaccAfter{} (six classes), but it removes veracity information with it. In these benchmarks corpus identity and label are entangled: each corpus collected its fake and real items from different outlets, periods and formats, so the directions that identify a corpus are partly the directions that separate its classes. Invariance to the corpus is therefore not invariance to the shortcut, and the remaining ranking gap is a property of the data rather than of the learner.

\begin{table}[t]
\centering
\caption{Leave-one-corpus-out repairs (\textsc{tfidf-lr} features; mean over the seven held-out veracity corpora and five seeds; s.d.\ across seeds). $\Delta$AUC versus pooled training with Holm-corrected Wilcoxon $p$ over the 35 (target, seed) pairs. For reference, the mean single-source transfer AUC to the same targets is 0.567, the best single source (chosen with target labels) 0.671, and the targets' own in-domain AUC 0.914.}\label{tab:repair}
\small
\begin{tabular}{@{}lccccc@{}}
\toprule
Training & AUC & AP & F1$_{0.5}$ & F1$^\star$ & $\Delta$AUC ($p$) \\
\midrule
pooled (6 corpora) & 0.660\,{\scriptsize$\pm$0.013} & 0.582 & 0.581 & 0.634 & -- \\
+ (corpus, class) balancing & 0.659\,{\scriptsize$\pm$0.013} & 0.583 & 0.582 & 0.632 & -0.001 (1.000) \\
+ balancing + masking & 0.661\,{\scriptsize$\pm$0.012} & 0.584 & 0.576 & 0.633 & +0.001 (1.000) \\
+ balancing + corpus-feature deletion & 0.636\,{\scriptsize$\pm$0.011} & 0.556 & 0.560 & 0.611 & -0.023 (0.006) \\
SVD-300 + balancing & 0.647\,{\scriptsize$\pm$0.009} & 0.566 & 0.583 & 0.624 & -0.013 (0.043) \\
SVD-300 + INLP & 0.605\,{\scriptsize$\pm$0.011} & 0.520 & 0.525 & 0.591 & -0.054 (0.001) \\
SVD-300 + CORAL & 0.621\,{\scriptsize$\pm$0.008} & 0.539 & 0.563 & 0.604 & -0.038 ($<$0.001) \\
\bottomrule
\end{tabular}
\end{table}

\begin{table}[h]
\centering
\caption{Leave-one-corpus-out AUC per held-out corpus (mean of five seeds), against single-source transfer and the in-domain ceiling.}\label{tab:repair-targets}
\small
\begin{tabular}{@{}lccccccc@{}}
\toprule
Held out & mean single & best single & pooled & + masking & INLP & CORAL & in-domain \\
\midrule
ISOT & 0.670 & 0.847 & 0.865 & 0.890 & 0.805 & 0.781 & 0.999 \\
WELFake & 0.540 & 0.690 & 0.648 & 0.657 & 0.643 & 0.643 & 0.989 \\
PolitiFact & 0.592 & 0.730 & 0.677 & 0.655 & 0.600 & 0.650 & 0.893 \\
GossipCop & 0.522 & 0.544 & 0.538 & 0.546 & 0.558 & 0.520 & 0.869 \\
LIAR & 0.530 & 0.552 & 0.528 & 0.528 & 0.502 & 0.514 & 0.698 \\
Twitter15 & 0.563 & 0.657 & 0.647 & 0.645 & 0.592 & 0.626 & 0.981 \\
Twitter16 & 0.552 & 0.678 & 0.715 & 0.703 & 0.536 & 0.616 & 0.967 \\
\bottomrule
\end{tabular}
\end{table}

